% slots for James to write into
% the red todo text shows in the pdf so nothing gets missed and the % lines below are notes that never print
% every table and figure here has to be referenced somewhere in the prose

\paragraph{Accuracy.}
\todo{Present first, do not interpret yet. Reference Table~\ref{tab:accuracy} (required summary table, all 11 configurations overall and by difficulty), Figure~\ref{fig:size} and Figure~\ref{fig:difficulty}.}
% key numbers
% coder ex 45.3 55.8 67.9 80.0 and awq 80.4
% qwen3 off 60.2 71.2 74.0 and on 71.8 79.4 78.9 with seed means 72.0 79.0 79.3 (sd 0.3 0.3 0.5)
% every configuration falls from easy to extra hard, coder 7b goes 91.5 to 53.0, nothing passes 58.4 on extra hard
% em is 15 to 44 points below ex everywhere

\begin{table}[htbp]
\centering
\caption{Accuracy of every Task 1 configuration, in percent, overall and by difficulty.}
\label{tab:accuracy}
\scriptsize
\setlength{\tabcolsep}{2.6pt}
\begin{tabular}{@{}l rrrrr rrrrr c r r@{}}
\toprule
 & \multicolumn{5}{c}{EM} & \multicolumn{5}{c}{EX} & EX all & EX gold & EX seed \\
\cmidrule(lr){2-6}\cmidrule(lr){7-11}
Configuration & Easy & Med. & Hard & Extra & All & Easy & Med. & Hard & Extra & All & 95\% CI & values & mean $\pm$ SD \\
\midrule
Coder 0.5B & 59.3 & 25.3 & 23.6 & 6.0 & 30.1 & 65.3 & 46.0 & 36.8 & 22.3 & 45.3 & [42.3, 48.3] & 46.5 &  \\
Coder 1.5B & 56.9 & 30.7 & 18.4 & 7.2 & 31.1 & 73.4 & 62.3 & 35.1 & 33.7 & 55.8 & [52.8, 58.8] & 58.4 &  \\
Coder 3B & 71.0 & 26.0 & 29.3 & 4.8 & 33.9 & 83.1 & 72.2 & 59.8 & 42.2 & 67.9 & [65.0, 70.7] & 70.2 &  \\
Coder 7B & 78.6 & 49.8 & 43.7 & 12.0 & 49.6 & 91.5 & 85.4 & 75.3 & 53.0 & 80.0 & [77.4, 82.3] & 82.9 &  \\
Coder 7B AWQ & 77.0 & 51.8 & 43.7 & 9.0 & 49.6 & 91.9 & 86.1 & 77.0 & 51.2 & 80.4 & [77.8, 82.7] & 83.1 &  \\
\midrule
Qwen3-1.7B off & 54.8 & 11.7 & 9.8 & 1.2 & 20.0 & 83.9 & 63.0 & 44.3 & 33.7 & 60.2 & [57.1, 63.1] & 62.2 &  \\
Qwen3-1.7B on & 81.9 & 38.8 & 23.0 & 10.2 & 41.9 & 88.3 & 80.0 & 52.9 & 44.6 & 71.8 & [68.9, 74.4] & 74.6 & 72.0 $\pm$ 0.3 \\
Qwen3-4B off & 66.1 & 32.7 & 17.8 & 6.6 & 34.0 & 85.9 & 76.9 & 58.0 & 47.6 & 71.2 & [68.3, 73.9] & 73.4 &  \\
Qwen3-4B on & 84.3 & 44.4 & 31.0 & 15.7 & 47.1 & 91.9 & 84.1 & 69.5 & 58.4 & 79.4 & [76.8, 81.8] & 82.3 & 79.0 $\pm$ 0.3 \\
Qwen3-8B off & 62.5 & 28.9 & 10.9 & 2.4 & 29.7 & 85.5 & 80.3 & 62.6 & 51.8 & 74.0 & [71.2, 76.6] & 77.3 &  \\
Qwen3-8B on & 85.1 & 42.6 & 27.6 & 9.6 & 45.0 & 89.5 & 83.9 & 71.3 & 57.8 & 78.9 & [76.3, 81.3] & 82.4 & 79.3 $\pm$ 0.5 \\
\bottomrule
\end{tabular}


\vspace{2pt}
\begin{minipage}{\linewidth}\scriptsize
\textit{Note.} Questions per level: 248 easy, 446 medium, 174 hard, 166 extra hard. Thinking rows are the seed 42 run, which every test uses; the last column averages all seeds. The EX interval is a Wilson 95\% interval. EX with gold values is the Spider paper's definition of EX and an upper bound on plain EX.
\end{minipage}
\end{table}

\begin{figure}[htbp]
\centering
\includegraphics[width=\linewidth]{size_curves.png}
\caption{EX and EM against parameter count. Thinking points are seed means, and their seed ranges are narrower than the markers. Differences are tested in Table~\ref{tab:mcnemar}, not by comparing points.}
\label{fig:size}
\end{figure}

\begin{figure}[htbp]
\centering
\includegraphics[width=\linewidth]{accuracy_by_difficulty.png}
\caption{EX by difficulty. Left: Qwen2.5-Coder by size with the 4 bit model. Right: Qwen3 by size with thinking off and on (thinking bars are seed means).}
\label{fig:difficulty}
\end{figure}

\paragraph{Planned comparisons.}
\todo{Present Table~\ref{tab:mcnemar}. Say under pessimistic EX every significant comparison stays significant.}
% coder steps all significant on ex (holm under 0.001) but em misses 0.5b to 1.5b (holm 1.00) and 1.5b to 3b (0.36)
% thinking +11.6 +8.2 +4.9 all holm under 0.001
% coder 7b vs 4b thinking -0.6 [-3.6, 2.3] holm 1.00 and awq +0.4 holm 1.00
% 4b to 8b no thinking ex +2.8 holm 0.097 but em -4.4 holm 0.007 so em reverses the ranking
% thinking 4b to 8b -0.5 holm 1.00

\begin{table}[htbp]
\centering
\caption{Planned paired comparisons on the seed 42 runs.}
\label{tab:mcnemar}
\scriptsize
\setlength{\tabcolsep}{3pt}
\resizebox{\linewidth}{!}{\begin{tabular}{@{}l r c r r r r@{}}
\toprule
 & \multicolumn{4}{c}{EX} & \multicolumn{2}{c}{EM} \\
\cmidrule(lr){2-5}\cmidrule(lr){6-7}
Comparison (first $\to$ second) & Difference [95\% CI] & Only first / only second & Holm $p$ & Pessimistic Holm $p$ & Difference [95\% CI] & Holm $p$ \\
\midrule
Coder 0.5B $\to$ 1.5B & +10.5 [6.8, 14.2] & 98 / 207 & $<$0.001 & $<$0.001 & +1.1 [$-$2.4, 4.5] & 1.000 \\
Coder 1.5B $\to$ 3B & +12.1 [8.7, 15.5] & 73 / 198 & $<$0.001 & $<$0.001 & +2.8 [$-$0.6, 6.4] & 0.355 \\
Coder 3B $\to$ 7B & +12.1 [9.0, 15.2] & 39 / 164 & $<$0.001 & $<$0.001 & +15.7 [12.1, 19.3] & $<$0.001 \\
Coder 1.5B $\to$ 7B & +24.2 [20.6, 28.1] & 43 / 293 & $<$0.001 & $<$0.001 & +18.5 [14.6, 22.1] & $<$0.001 \\
Qwen3-1.7B off $\to$ on & +11.6 [8.6, 14.7] & 46 / 166 & $<$0.001 & $<$0.001 & +21.9 [18.4, 25.4] & $<$0.001 \\
Qwen3-4B off $\to$ on & +8.2 [5.4, 11.4] & 52 / 137 & $<$0.001 & $<$0.001 & +13.1 [10.0, 16.2] & $<$0.001 \\
Qwen3-8B off $\to$ on & +4.9 [2.2, 7.8] & 56 / 107 & $<$0.001 & $<$0.001 & +15.3 [12.1, 18.6] & $<$0.001 \\
Qwen3 off, 1.7B $\to$ 4B & +11.0 [7.5, 14.5] & 72 / 186 & $<$0.001 & $<$0.001 & +14.0 [10.6, 17.4] & $<$0.001 \\
Qwen3 off, 4B $\to$ 8B & +2.8 [0.2, 5.4] & 63 / 92 & 0.097 & 0.220 & $-$4.4 [$-$7.3, $-$1.5] & 0.007 \\
Qwen3 on, 1.7B $\to$ 4B & +7.6 [4.8, 10.5] & 45 / 124 & $<$0.001 & $<$0.001 & +5.2 [2.2, 8.1] & 0.001 \\
Qwen3 on, 4B $\to$ 8B & $-$0.5 [$-$2.4, 1.3] & 48 / 43 & 1.000 & 1.000 & $-$2.1 [$-$4.7, 0.3] & 0.355 \\
Coder 7B $\to$ Qwen3-4B on & $-$0.6 [$-$3.6, 2.3] & 83 / 77 & 1.000 & 1.000 & $-$2.5 [$-$6.6, 1.3] & 0.375 \\
Coder 7B bf16 $\to$ AWQ & +0.4 [$-$0.7, 1.5] & 14 / 18 & 1.000 & 1.000 & +0.0 [$-$1.6, 1.6] & 1.000 \\
\bottomrule
\end{tabular}
}

\vspace{2pt}
\begin{minipage}{\linewidth}\scriptsize
\textit{Note.} Differences are second minus first, in points, with a 95\% bootstrap interval that resamples the 551 gold queries. ``Only first / only second'' counts the questions only one run got right, which are the only questions the exact McNemar test uses. Holm correction is applied separately to the 13 EX and the 13 EM tests. The pessimistic Holm $p$ column repeats the EX test with pessimistic EX, where every EX pass on a trivial gold answer that EM rejects is counted as wrong.
\end{minipage}
\end{table}

\paragraph{Efficiency.}
\todo{Present Table~\ref{tab:efficiency} and Figure~\ref{fig:cost}. Times are amortized over batches of 100 (see Inference settings).}
% 4b thinking 2.11 s and 867 tokens vs 0.055 s and 38 tokens off, about 38 times the time and 23 times the tokens
% coder 7b gets 80.0 with 34 tokens, 26 times fewer than 4b thinking
% awq 0.041 s vs 0.078 s bf16, 1.9 times faster, same accuracy
% coder 1.5b writes about 230 tokens since it explains itself, which is why its slower than the 3b

\begin{table}[htbp]
\centering
\caption{Generation cost per configuration on one NVIDIA L4.}
\label{tab:efficiency}
\small
\begin{tabular}{@{}l rrrr@{}}
\toprule
  & Seconds per & \multicolumn{2}{c}{Output tokens per question} & Reached \\ Seconds per & \multicolumn{2}{c}{Output tokens per question} & Reached \\
\cmidrule(lr){3-4}
Configuration & question & Mean & Median & token limit \\
\midrule
Coder 0.5B & 0.015 & 40 & 31 & 8 \\
Coder 1.5B & 0.088 & 230 & 253 & 6 \\
Coder 3B & 0.035 & 38 & 35 & 0 \\
Coder 7B & 0.078 & 34 & 28 & 0 \\
Coder 7B AWQ & 0.041 & 34 & 27.5 & 0 \\
\midrule
Qwen3-1.7B off & 0.033 & 40 & 35 & 1 \\
Qwen3-1.7B on & 1.27 & 874 & 571.5 & 2 \\
Qwen3-4B off & 0.055 & 38 & 33 & 0 \\
Qwen3-4B on & 2.11 & 867 & 540 & 0 \\
Qwen3-8B off & 0.092 & 38 & 34 & 0 \\
Qwen3-8B on & 5.46 & 1,045 & 644 & 2 \\
\bottomrule
\end{tabular}


\vspace{2pt}
\begin{minipage}{0.9\linewidth}\footnotesize
\textit{Note.} The seconds per question is calculated from batch time divided by batch size and measures throughput. Mean output tokens include the reasoning trace for thinking configurations. Reached token limit counts the outputs that were cut off at the output limit, which is 512 tokens without thinking and 8,192 tokens with thinking.
\end{minipage}
\end{table}

\begin{figure}[htbp]
\centering
\includegraphics[width=0.75\linewidth]{cost_tradeoff.png}
\caption{EX against mean output tokens per question (log scale). Thinking points are seed means.}
\label{fig:cost}
\end{figure}

\paragraph{Return on thinking.}
\todo{Present Table~\ref{tab:thinking}.}
% overall gain positive at every size, largest on hard (+10 to +13)
% per difficulty intervals are wide so describe it as a pattern only
% per 1,000 extra tokens extra hard gives the lowest return at every size

\begin{table}[htbp]
\centering
\caption{What turning thinking on buys at each Qwen3 size, by difficulty.}
\label{tab:thinking}
\scriptsize
\setlength{\tabcolsep}{3pt}
\begin{tabular}{@{}l ccccc@{}}
\toprule
Size & Easy & Medium & Hard & Extra & All \\
\midrule
\multicolumn{6}{@{}l}{\textit{EX gain from thinking, points [95\% CI]}} \\
1.7B & +3.5 [$-$1.0, 8.0] & +17.3 [13.0, 21.8] & +13.2 [5.9, 20.7] & +8.2 [0.4, 16.1] & +11.9 [9.0, 14.8] \\
4B & +5.2 [1.0, 9.6] & +6.5 [2.8, 10.4] & +12.8 [2.9, 22.0] & +9.6 [0.7, 18.3] & +7.8 [5.0, 10.7] \\
8B & +5.0 [1.1, 9.3] & +3.2 [$-$0.2, 6.7] & +10.3 [1.9, 18.6] & +6.2 [$-$1.0, 13.8] & +5.3 [2.7, 8.0] \\
\midrule
\multicolumn{6}{@{}l}{\textit{Extra output tokens per question}} \\
1.7B & 474 & 613 & 1,234 & 1,498 & 826 \\
4B & 430 & 655 & 1,264 & 1,479 & 836 \\
8B & 508 & 814 & 1,420 & 1,758 & 994 \\
\midrule
\multicolumn{6}{@{}l}{\textit{EX points gained per 1,000 extra tokens}} \\
1.7B & 7.4 & 28.3 & 10.7 & 5.5 & 14.4 \\
4B & 12.2 & 10.0 & 10.1 & 6.5 & 9.3 \\
8B & 9.8 & 3.9 & 7.3 & 3.5 & 5.3 \\
\bottomrule
\end{tabular}


\vspace{2pt}
\begin{minipage}{\linewidth}\scriptsize
\textit{Note.} Thinking accuracy is averaged over every seed of a size for each question before comparing. Intervals resample the 551 gold queries. Each difficulty level has only 166 to 446 questions, so the per level values are descriptive.
\end{minipage}
\end{table}

\paragraph{EM against EX.}
\todo{Present Table~\ref{tab:emparser} in two sentences.}
% the 2018 parser cant read 32 to 65 percent of predictions
% for qwen3 8b no thinking 558 of 577 unreadable queries run fine so its the parser and not broken sql
% qwen3 without thinking loses about half its correct answers to the parser (55.5 percent at 1.7b, 50.7 at 8b)



\begin{table}[htbp]
\centering
\caption{Predictions Spider's 2018 SQL parser cannot read, which EM scores as wrong.}
\label{tab:emparser}
\small
\begin{tabular}{@{}l rrrrr@{}}
\toprule
 & & Cannot be & Of those, & Of those, & EX-correct answers \\
Configuration & Predictions & parsed (\%) & run & EX right & lost to the parser (\%) \\
\midrule
Coder 0.5B & 1034 & 507 (49.0) & 124 & 65 & 13.9 \\
Coder 1.5B & 1034 & 516 (49.9) & 293 & 158 & 27.4 \\
Coder 3B & 1034 & 499 (48.3) & 379 & 246 & 35.0 \\
Coder 7B & 1034 & 328 (31.7) & 291 & 208 & 25.2 \\
Coder 7B AWQ & 1034 & 336 (32.5) & 305 & 214 & 25.8 \\
\midrule
Qwen3-1.7B off & 1034 & 667 (64.5) & 543 & 345 & 55.5 \\
Qwen3-1.7B on & 1032 & 418 (40.5) & 368 & 250 & 33.7 \\
Qwen3-4B off & 1034 & 451 (43.6) & 410 & 262 & 35.6 \\
Qwen3-4B on & 1034 & 403 (39.0) & 395 & 267 & 32.5 \\
Qwen3-8B off & 1034 & 577 (55.8) & 558 & 388 & 50.7 \\
Qwen3-8B on & 1032 & 443 (42.9) & 441 & 308 & 37.7 \\
\bottomrule
\end{tabular}


\vspace{2pt}
\begin{minipage}{0.9\linewidth}\footnotesize
\textit{Note.} ``Lost to the parser'' is the share of EX-correct answers that EM scores as wrong only because the query cannot be parsed. Predictions excludes outputs with no prediction.
\end{minipage}
\end{table}

\paragraph{Errors and robustness.}
\todo{Present Table~\ref{tab:errors}.}
% crash share of wrong answers falls from 68.9 (coder 0.5b) to 4.2 (4b thinking) and 0.9 (8b thinking)
% 81.7 percent of 4b thinking errors are structural, gold values add only 1 to 3 points
% pessimistic ex drops 2.0 to 5.4 points and changes no conclusion

\begin{table}[htbp]
\centering
\caption{Wrong answers by cause, and EX when possibly accidental passes are counted as wrong.}
\label{tab:errors}
\small
\setlength{\tabcolsep}{4pt}
\begin{tabular}{@{}l rrr rrr rr@{}}
\toprule
 & & No & & \multicolumn{3}{c}{Share of wrong answers (\%)} & Trivial gold, & Pessimistic \\
\cmidrule(lr){5-7}
Configuration & Wrong & prediction & Timeout & Crash & Value only & Structural & EM rejects & EX \\
\midrule
Coder 0.5B & 566 & 0 & 0 & 68.9 & 2.3 & 28.8 & 35 & 41.9 \\
Coder 1.5B & 457 & 0 & 0 & 49.7 & 5.9 & 44.4 & 43 & 51.6 \\
Coder 3B & 332 & 0 & 0 & 37.0 & 7.2 & 55.7 & 46 & 63.4 \\
Coder 7B & 207 & 0 & 0 & 18.4 & 14.5 & 67.1 & 23 & 77.8 \\
Coder 7B AWQ & 203 & 0 & 0 & 15.8 & 13.8 & 70.4 & 23 & 78.1 \\
\midrule
Qwen3-1.7B off & 412 & 0 & 0 & 31.1 & 5.1 & 63.8 & 56 & 54.7 \\
Qwen3-1.7B on & 292 & 2 & 1 & 19.9 & 9.6 & 69.5 & 40 & 67.9 \\
Qwen3-4B off & 298 & 0 & 0 & 14.8 & 7.7 & 77.5 & 39 & 67.4 \\
Qwen3-4B on & 213 & 0 & 0 & 4.2 & 14.1 & 81.7 & 26 & 76.9 \\
Qwen3-8B off & 269 & 0 & 0 & 7.8 & 12.6 & 79.6 & 43 & 69.8 \\
Qwen3-8B on & 218 & 2 & 0 & 0.9 & 16.5 & 81.7 & 28 & 76.2 \\
\bottomrule
\end{tabular}


\vspace{2pt}
\begin{minipage}{0.95\linewidth}\footnotesize
\textit{Note.} Each wrong answer is in exactly one category, checked in order: no prediction, timeout, crash (the query errors), value only (correct once gold values are substituted), structural. 73 of the 1,034 gold answers are empty or a single 0 or NULL. Pessimistic EX counts every EX pass on those that EM also rejects as wrong, so it is a bound and not an estimate.
\end{minipage}
\end{table}

\paragraph{Answering RQ1.}
\todo{Interpret, with numbers from the tables. One claim per point: top tier is a statistical tie; size helps steadily for Coder; thinking helps at every size but less as size grows, worth about one size step; difficulty; cost; EM is unreliable for these models; what is left is mostly logic errors; robustness; then a short limitations paragraph (single database EX, possible pretraining contamination, one prompt without sample rows, one family per model type, no MoE, dev split only).}
% one sentence answer from the outline
% the best open models at 8b or fewer reach about 80 percent ex zero shot on unseen databases by two routes at very different cost, and em understates all of them
% zhong et al 2020 found only a 2.5 percent em false negative rate for 2020 fine tuned models which copied spiders gold style but zero shot llms dont
